%Revised version
\documentclass[12pt]{amsart}
\usepackage{latexsym,enumerate}
\usepackage{amssymb, xcolor,mathrsfs}
\usepackage{amsmath,amsthm,amsfonts,amssymb,latexsym, mathabx}
%\usepackage{amsfonts}
%\usepackage{showlabels}
\usepackage[pagebackref,colorlinks]{hyperref}
%\usepackage[T1]{fontenc}
%\usepackage[utf8]{inputenc}
%\usepackage{mathtools}
\usepackage{arydshln}

\headheight=7pt \textheight=574pt \textwidth=432pt \topmargin=14pt
\oddsidemargin=18pt \evensidemargin=18pt

\newtheorem{theorem}{Theorem}[section]
\newtheorem{lemma}[theorem]{Lemma}
\newtheorem{proposition}[theorem]{Proposition}
\newtheorem{corollary}[theorem]{Corollary}

\newtheorem{theorema}{Theorem}
\renewcommand{\thetheorema}{\Alph{theorema}}
\newtheorem{conj}[theorema]{Conjecture}

\newtheorem{consequence}[theorem]{Consequence}
\newtheorem{conjecture}[theorem]{Conjecture}
\newtheorem{problem}[theorem]{Problem}
\newtheorem{question}[theorem]{Question}
\newtheorem{MainTheorem}{Theorem}
\renewcommand{\theMainTheorem}{\Alph{MainTheorem}}

\theoremstyle{definition}
\newtheorem{definition}[theorem]{Definition}
\newtheorem{example}[theorem]{Example}
\newtheorem{exercise}[theorem]{Exercise}
\newtheorem{notation}[theorem]{Notation}
\newtheorem{remark}[theorem]{Remark}
\newtheorem{observation}[theorem]{Observation}
\newtheorem{hypothesis}[theorem]{Hypothesis}


\numberwithin{equation}{section}

\renewcommand{\theequation}{\arabic{section}.\arabic{equation}}
\newcommand{\GL}{{\mathrm {GL}}}
\newcommand{\PGL}{{\mathrm {PGL}}}
\newcommand{\SL}{{\mathrm {SL}}}
\newcommand{\PSL}{{\mathrm {PSL}}}
\newcommand{\POmega}{{\mathrm {P\Omega}}}


\newcommand{\GU}{{\mathrm {GU}}}
\newcommand{\PGU}{{\mathrm {PGU}}}
\newcommand{\SU}{{\mathrm {SU}}}
\newcommand{\PSU}{{\mathrm {PSU}}}
\newcommand{\U}{{\mathrm {U}}}

\newcommand{\GO}{{\mathrm {GO}}}
\newcommand{\Spin}{{\mathrm {Spin}}}
\newcommand{\SO}{{\mathrm {SO}}}
\newcommand{\PCO}{{\mathrm {PCO}}}
\newcommand{\CB}{\mathbf{C}}


\newcommand{\Sp}{{\mathrm {Sp}}}
\newcommand{\PSp}{{\mathrm {PSp}}}
\newcommand{\PCSp}{{\mathrm {PCSp}}}

\newcommand{\rk}{{\mathrm {rk}}}
\newcommand{\Ker}{\operatorname{Ker}}
\newcommand{\Aut}{{\mathrm {Aut}}}
\newcommand{\Out}{{\mathrm {Out}}}
\newcommand{\Mult}{{\mathrm {Mult}}}
%\newcommand{\Stab}{{\mathrm {Stab}}}
\newcommand{\Schur}{{\mathrm {Schur}}}
\newcommand{\Inn}{{\mathrm {Inn}}}
\newcommand{\Irr}{{\mathrm {Irr}}}
\newcommand{\IBR}{{\mathrm {IBr}}}
\newcommand{\IBRL}{{\mathrm {IBr}}_{\ell}}
\newcommand{\ord}{{\mathrm {ord}}}
\def\id{\mathop{\mathrm{ id}}\nolimits}
\renewcommand{\Im}{{\mathrm {Im}}}
\newcommand{\Ind}{{\mathrm {Ind}}}
\newcommand{\diag}{{\mathrm {diag}}}
\newcommand{\soc}{{\mathrm {soc}}}
\newcommand{\Soc}{{\mathrm {Soc}}}
\newcommand{\End}{{\mathrm {End}}}
\newcommand{\sol}{{\mathrm {sol}}}
\newcommand{\Hom}{{\mathrm {Hom}}}
\newcommand{\Mor}{{\mathrm {Mor}}}
\def\rank{\mathop{\mathrm{ rank}}\nolimits}
\newcommand{\Syl}{{\mathrm {Syl}}}
\newcommand{\St}{{\mathsf {St}}}
\newcommand{\Stab}{{\mathrm {Stab}}}
\newcommand{\Tr}{{\mathrm {Tr}}}
\newcommand{\tr}{{\mathrm {tr}}}
\newcommand{\cl}{{\mathrm {cl}}}
\newcommand{\Cl}{{\mathrm {Cl}}}
\newcommand{\cd}{{\mathrm {cd}}}
\newcommand{\lcm}{{\mathrm {lcm}}}
\newcommand{\Gal}{{\rm Gal}}
\newcommand{\Spec}{{\mathrm {Spec}}}
\newcommand{\ad}{{\mathrm {ad}}}
\newcommand{\Sym}{{\mathrm {Sym}}}
\newcommand{\Alt}{{\mathrm {Alt}}}
\newcommand{\Char}{{\mathrm {Char}}}
\newcommand{\Lin}{{\mathrm {Lin}}}
\newcommand{\sign}{{\mathrm {sgn}}}

\newcommand{\pr}{{\mathrm {pr}}}
\newcommand{\rad}{{\mathrm {rad}}}
\newcommand{\abel}{{\mathrm {abel}}}
\newcommand{\codim}{{\mathrm {codim}}}
\newcommand{\ind}{{\mathrm {ind}}}
\newcommand{\Res}{{\mathrm {Res}}}
\newcommand{\Ann}{{\mathrm {Ann}}}
\newcommand{\CC}{{\mathbb C}}
\newcommand{\RR}{{\mathbb R}}
\newcommand{\QQ}{{\mathbb Q}}
\newcommand{\ZZ}{{\mathbb Z}}
\newcommand{\NN}{{\mathbb N}}
\newcommand{\EE}{{\mathbb E}}
\newcommand{\PP}{{\mathbb P}}
\newcommand{\GG}{{\mathbb G}}
\newcommand{\HH}{{\mathbb HH}}

\newcommand{\SSS}{{\mathbb S}}
\newcommand{\AAA}{{\mathbb A}}
\newcommand{\HA}{\hat{{\mathbb A}}}
\newcommand{\FF}{{\mathbb F}}
\newcommand{\FP}{\mathbb{F}_{p}}
\newcommand{\FQ}{\mathbb{F}_{q}}
\newcommand{\FQB}{\mathbb{F}_{q}^{\bullet}}
\newcommand{\FQA}{\mathbb{F}_{q^{2}}}
\newcommand{\KK}{{\mathbb K}}
\newcommand{\EC}{\mathcal{E}}
\newcommand{\XC}{\mathcal{G}}
\newcommand{\HC}{\mathcal{H}}
\newcommand{\AC}{\mathcal{A}}
\newcommand{\OC}{\mathcal{O}}
\newcommand{\PC}{\mathcal{P}}

\newcommand{\OL}{O^{\ell'}}
\newcommand{\IL}{\mathcal{I}}
\newcommand{\CL}{\mathcal{C}}
\newcommand{\RL}{\mathcal{R}}
\newcommand{\TL}{\mathcal{T}}
\newcommand{\TC}{\mathcal{T}}
\newcommand{\GC}{\mathcal{G}}
%\newcommand{\GCR}{\mathcal{R}}
\newcommand{\eps}{\epsilon}
\newcommand{\LC}{\mathcal{L}}
\newcommand{\M}{\mathcal{M}}
%\newcommand{\HC}{\mathcal{H}}
%\newcommand{\bL}{\bar{L}}
\newcommand{\Om}{\Omega}
\newcommand{\tn}{\hspace{0.5mm}^{t}\hspace*{-1mm}}
\newcommand{\ta}{\hspace{0.5mm}^{2}\hspace*{-0.2mm}}
\newcommand{\tb}{\hspace{0.5mm}^{3}\hspace*{-0.2mm}}
\newcommand{\p}{^{\prime}}
\newcommand{\Center}{\mathbf{Z}}
\newcommand{\Centralizer}{\mathbf{C}}
\newcommand{\Normalizer}{\mathbf{N}}
\newcommand{\acd}{\mathsf{acd}}
\newcommand{\acde}{\mathsf{acd}_{\mathrm {even}}}
\newcommand{\acdo}{\mathsf{acd}_{\mathrm {odd}}}
\newcommand{\odd}{\mathrm{odd}}

\newcommand{\wt}{\widetilde}

\newcommand{\height}{\mathbf{ht}}
\newcommand{\lev}{\mathrm{\mathbf{lev}}}
\newcommand{\Bl}{\mathrm{\mathrm{Bl}}}


\newcommand{\acs}{\mathrm{acs}}
\newcommand{\bC}{{\mathbf{C}}}
\newcommand{\bG}{{\mathbf{G}}}
\newcommand{\bS}{{\mathbf{S}}}
\newcommand{\bH}{{\mathbf{H}}}
\newcommand{\bT}{{\mathbf{T}}}
\newcommand{\bL}{{\mathbf{L}}}
\newcommand{\bO}{{\mathbf{O}}}
\newcommand{\bN}{{\mathbf{N}}}
\newcommand{\bB}{{\mathbf{B}}}
\newcommand{\bU}{{\mathbf{U}}}
\newcommand{\NB}{{\mathbf{N}}}
\newcommand{\bZ}{{\mathbf{Z}}}
\newcommand{\bX}{{\mathbf{X}}}
\newcommand{\Al}{\textup{\textsf{A}}}
\newcommand{\Sy}{\textup{\textsf{S}}}
\newcommand{\tw}[1]{{}^#1}
\renewcommand{\mod}{\bmod \,}

\newcommand{\bff}{\mathbf{f}}

\newcommand{\pcore}{\mathbf{O}}
\newcommand{\GammaL}{\operatorname{\Gamma L}}
\newcommand{\IBr}{\operatorname{IBr}}
\newcommand{\bga}{\boldsymbol{\gamma}}
\newcommand{\blam}{\boldsymbol{\lambda}}
\newcommand{\lam}{\lambda}
\makeatletter \@namedef{subjclassname@2020}{\textup{2020}
Mathematics Subject Classification} \makeatother

\def\nor{\trianglelefteq\,}

\newcommand{\hungcomment}{\textcolor{red}}


\title[Hall $\pi$-subgroups and characters of $\pi'$-degree]{Hall $\pi$%%
%% Correction 0 [ ]
%% Annotated text: "HALL π-SUBGROUPS AND<Caret> CHARACTERS OF"
%% Comment: "" 
%% 
%% START of correction 0
-subgroups and characters %%
%% END of correction 0
of $\pi'$-degree}

\author[Eugenio Giannelli]{Eugenio Giannelli}
\address{Dipartimento di Matematica e Informatica U. Dini, Firenze,
Italy}
\email{eugenio.giannelli@unifi.it}

\author[Nguyen N. Hung]{Nguyen N. Hung}
\address{Department of Mathematics, The University of Akron, Akron,
OH 44325, USA}
\email{hungnguyen@uakron.edu}


\thanks{The first author's research is funded by: the European Union Next
Generation EU, M4C1, CUP B53D23009410006, PRIN 2022 - 2022PSTWLB
Group Theory and Applications; and the INDAM-GNSAGA Project CUP
E53C24001950001. The second author gratefully acknowledges support
from the AMS--Simons Research Enhancement Grant (AWD-000167 AMS). We
thank Gunter Malle for several helpful comments on an earlier
version, particularly regarding
Proposition~\ref{prop:linear-unitary}}



\subjclass[2020]{Primary 20C15, 20C30, 20C33, 20D20}
%
\keywords{Characters, fields of values, Hall $\pi$-subgroups,
characters of $\pi'$-degree.}

%\date{}

\begin{document}
\begin{abstract}
%%
%% Correction 1 [ ]
%% Annotated text: "Abstract. We study<Caret> the relationship"
%% Comment: "" 
%% 
%% START of correction 1
We study the %%
%% END of correction 1
relationship between the existence of Hall
$\pi$-subgroups and that of irreducible characters %%
%% Correction 3 [ ]
%% Annotated text: "characters of π1-<Caret>degree with prescribed"
%% Comment: "" 
%% 
%% Correction 4 [ ]
%% Annotated text: "characters of π<Caret>1-degree with prescribed"
%% Comment: "" 
%% 
%% Correction 5 [ ]
%% Annotated text: "π1-degree with p<Caret>rescribed fields of"
%% Comment: "" 
%% 
%% START of corrections 3, 4, 5
of $\pi'$-degree
with prescribed fields %%
%% END of corrections 3, 4, 5
of values in finite groups. This %%
%% Correction 2 [ ]
%% Annotated text: "This work extends<Caret> a result"
%% Comment: "" 
%% 
%% START of correction 2
work extends
a %%
%% END of correction 2
result of Navarro and Tiep from a single odd prime to multiple odd
primes.
\end{abstract}

\maketitle

%\tableofcontents

%%%%%%%%%%%%%%%%%%%%%%%%
%%%%%%%%%%%%%%%%%%%%%%%

\section{Introduction}

%%
%% Correction 6 [ ]
%% Annotated text: "A w<Caret>ell-known conjecture of"
%% Comment: "" 
%% 
%% Correction 7 [ ]
%% Annotated text: "A well<Caret>-known conjecture of"
%% Comment: "" 
%% 
%% Correction 8 [ ]
%% Annotated text: "A well-known conje<Caret>cture of R."
%% Comment: "" 
%% 
%% START of corrections 6, 7, 8
A well-known conjecture of %%
%% END of corrections 6, 7, 8
R.~Gow, proved by G.~Navarro and
P.~H.~Tiep~\cite{Navarro-Tiep2}, asserts that if $G$ is a finite
group of even order, then $G$ possesses a nontrivial (complex)
irreducible character of odd degree whose values are rational. This
result was subsequently generalized to all primes by the same
authors, as %%
%% Correction 9 [ ]
%% Annotated text: "Theorem 1.<Caret>1 ([NT06], Theorem"
%% Comment: "" 
%% 
%% Correction 10 [ ]
%% Annotated text: "Theorem 1.1 ([NT06],<Caret> Theorem A)."
%% Comment: "" 
%% 
%% START of corrections 9, 10
follows.

\begin{theorem}[\cite{Navarro-Tiep1}, Theorem~A]\label{thm:Navarro-Tiep}
Let $G$ be a finite group of order divisible by a prime %%
%% END of corrections 9, 10
$p$. Then
$G$ has a nontrivial irreducible character of degree not divisible
by $p$ whose values lie %%
%% Correction 11 [ ]
%% Annotated text: "lie in Qpe2πi{pq<Caret>."
%% Comment: "" 
%% 
%% Correction 12 [ ]
%% Annotated text: "lie in Q<Caret>pe2πi{pq."
%% Comment: "" 
%% 
%% START of corrections 11, 12
in $\QQ(e^{2\pi i/p})$.
\end{theorem}

In%%
%% END of corrections 11, 12
~\cite{GHSV21}, together with A.\,A.~Schaeffer Fry and C.~Vallejo,
we attempted to extend this theorem from a single prime $p$ to a set
$\pi$ of two primes. Among %%
%% Correction 13 [ ]
%% Annotated text: "π “ t2,<Caret>p and gcdp|G|,"
%% Comment: "" 
%% 
%% Correction 14 [ ]
%% Annotated text: "“ t2, pu<Caret> and gcdp|G|,"
%% Comment: "" 
%% 
%% Correction 15 [ ]
%% Annotated text: "2pq ą 1<Caret>, then G"
%% Comment: "" 
%% 
%% Correction 16 [ ]
%% Annotated text: "2pq ą 1,<Caret> then G"
%% Comment: "" 
%% 
%% Correction 17 [ ]
%% Annotated text: "1, then G<Caret> possesses a"
%% Comment: "" 
%% 
%% Correction 18 [ ]
%% Annotated text: "ą 1, then<Caret> G possesses"
%% Comment: "" 
%% 
%% START of corrections 13, 14, 15, 16, 17, 18
other results, it was shown that if
$\pi=\{2,p\}$ and $\gcd(|G|,2p)>1$, then $G$ possesses %%
%% END of corrections 13, 14, 15, 16, 17, 18
a nontrivial
irreducible character of $\pi'$-degree whose values are contained %%
%% Correction 19 [ ]
%% Annotated text: "contained in Qpe2πi{pq<Caret>. (Here, a"
%% Comment: "" 
%% 
%% Correction 20 [ ]
%% Annotated text: "contained in Qpe2πi{pq.<Caret> (Here, a"
%% Comment: "" 
%% 
%% Correction 21 [ ]
%% Annotated text: "contained in Q<Caret>pe2πi{pq. (Here, a"
%% Comment: "" 
%% 
%% START of corrections 19, 20, 21
in
$\QQ(e^{2\pi i/p})$. (Here, %%
%% END of corrections 19, 20, 21
a character $\chi$ is said to have
$\pi'$-degree if $\chi(1)$ is not divisible by any prime in $\pi$.)
Unfortunately, this phenomenon does not hold in general when $\pi$
consists of two odd primes. For example, as pointed out in
\cite[Proposition~4.1]{GHSV21}, the Tits group ${}^2F_4(2)'$ has no
nontrivial irreducible character of %%
%% Correction 26 [ ]
%% Annotated text: "of t3, 5u1-<Caret>degree with values"
%% Comment: "" 
%% 
%% Correction 27 [ ]
%% Annotated text: "of t3, 5u1-degree<Caret> with values"
%% Comment: "" 
%% 
%% START of corrections 26, 27
$\{3,5\}'$-degree with %%
%% END of corrections 26, 27
values in
$\QQ(e^{2\pi i/15})$.

In this paper, we propose a different extension -- perhaps a more
natural one -- of Theorem~\ref{thm:Navarro-Tiep} that works for
an arbitrary set of odd primes. In the following, we write $\pi(G)$
for the set of prime divisors of $|G|$. Recall that a subgroup $H\le
G$ is called %%
%% Correction 28 [ ]
%% Annotated text: "called a Hall<Caret> π-subgroup of"
%% Comment: "" 
%% 
%% Correction 29 [ ]
%% Annotated text: "a Hall π-subgro<Caret>up of G"
%% Comment: "" 
%% 
%% START of corrections 28, 29
a Hall $\pi$-subgroup of %%
%% END of corrections 28, 29
$G$ %%
%% Correction 22 [ ]
%% Annotated text: "G if πp<Caret>Hq Ď π"
%% Comment: "" 
%% 
%% Correction 23 [ ]
%% Annotated text: "G if πpH<Caret>q Ď π"
%% Comment: "" 
%% 
%% Correction 24 [ ]
%% Annotated text: "π and |G<Caret> H| is"
%% Comment: "" 
%% 
%% Correction 25 [ ]
%% Annotated text: "and |G :<Caret> is not"
%% Comment: "" 
%% 
%% START of corrections 22, 23, 24, 25
if $\pi(H)\subseteq \pi$
and $|G:H|$ is %%
%% END of corrections 22, 23, 24, 25
not divisible by any prime in $\pi$.

\begin{theorema}\label{thm:main}
Let $\pi$ be a set of odd primes and $G$ a finite group possessing a
nontrivial Hall $\pi$-subgroup. Then $G$ has a nontrivial
$\pi'$-degree irreducible character with values in $\QQ(e^{2\pi
i/p})$ for some $p\in\pi$.
\end{theorema}



Theorem~\ref{thm:main} may be viewed as a generalization of
Theorem~\ref{thm:Navarro-Tiep}. In fact, when $p$ is odd,
Theorem~\ref{thm:Navarro-Tiep} follows from the case $\pi=\{p\}$ of
Theorem~\ref{thm:main} together with the first Sylow theorem.

If we assume that $2\in \pi$, then the conclusion of
Theorem~\ref{thm:main} remains true when %%
%% Correction 30 [ ]
%% Annotated text: "|π| ď 2,<Caret> as noted"
%% Comment: "" 
%% 
%% START of correction 30
$|\pi|\leq 2$, as %%
%% END of correction 30
noted
above, but fails in general once $|\pi|\geq 3$. For example, if $G$
is any non-abelian simple group and %%
%% Correction 31 [ ]
%% Annotated text: "π “ πpGq,<Caret> then the"
%% Comment: "" 
%% 
%% START of correction 31
$\pi=\pi(G)$, then %%
%% END of correction 31
the trivial
character $1_G$ is the only irreducible character of $G$ with
$\pi'$-degree. There are also counterexamples with $\pi\varsubsetneq
\pi(G)$, such as $(G,\pi)=(\Al_7,\{2,3,5\})$. At present, we do not
have a conceptual explanation for why the prime $2$ is special in
this context.

Similar to Gow's conjecture and the results in
\cite{Navarro-Tiep1,GHSV21}, Theorem~\ref{thm:main} admits a clean
reduction to finite simple groups. Accordingly, the main body of
this paper is devoted to proving the result for such groups. This
requires some new work on the relationship between the existence of
Hall $\pi$-subgroups and $\pi'$-degree irreducible characters in
simple groups of Lie type.


%%%%%%%%%%%%%%%%%%%%%%%
%%%%%%%%%%%%%%%%%%%%%%%
%\section{Alternating groups}
%For notational simplicity, we write $\QQ(\zeta_k)$ for the $k$-th
%cyclotomic field $\QQ(e^{2\pi i/k})$. As usual, $\Irr(G)$ denotes
%the set of all irreducible characters of $G$, and $\Irr_{\pi'}(G)$
%denotes the subset consisting of characters whose degrees are not
%divisible by any prime in $\pi$. Given $x,y\in\mathbb{Z}$ we let $[x,y]:=\{z\in\mathbb{Z}\ |\ x\leq z\leq y\}$.
%
%\begin{lemma}\label{lem:noHall}
%Let $n\in\mathbb{N}_{\geq 5}$ and let $\pi\subseteq [1,n]$ be a set of prime numbers. Assume that $2\notin\pi$.
%Then the alternating group $\Al_n$ admits a Hall $\pi$-subgroup if and only if $|\pi|=1$.
%\end{lemma}
%\begin{proof}
%If $|\pi|=1$ the statement is obviously implied by Sylow theory.
%Let us now assume that $\Al_n$ admits a Hall $\pi$-subgroup $H$ and let us suppose for a contradiction that $|\pi|\geq 2$.
%In particular let $\pi=\{p_1,p_2,\ldots, p_t\}$ for some odd primes $p_1<p_2<\cdots<p_t$.
%We proceed by distinguishing two cases, depending on the solvability of $H$.
%
%\smallskip
%
%%\noindent\textit{Case} (i).
%\noindent $\bullet$ Let us first assume that $H$ is a solvable group. By \cite{Hall1} we know that $H$ admits a $\{p_1,p_2\}$-Hall subgroup $K$. It follows that $K$ is a solvable $\{p_1,p_2\}$-Hall subgroup of $\Al_n$. Moreover, since $2\notin \{p_1,p_2\}$ we also have that $K$ is a solvable $\{p_1,p_2\}$-Hall subgroup of $\Sy_n$. Using \cite[Theorem A4]{Hall2}, we deduce that $|K|$ is even and therefore that $2\in \{p_1,p_2\}\subseteq \pi$. This clearly contradicts our hypothesis.
%
%
%
%\smallskip
%
%\noindent $\bullet$ Let us now consider the case where $H$ is a non-solvable Hall $\pi$-subgroup of $\Al_n$. Again, since $2\notin \pi$ we have that $H$ is a non-solvable Hall $\pi$-subgroup of $\Sy_n$.
%This family of subgroups of symmetric groups is completely described in \cite{Thompson}. In particular, we deduce that either $H=\Sy_n$ or $n$ is a prime number and $H=\Sy_{n-1}$. In both cases we would have that $2\in\pi$, contradicting our hypothesis.
%\end{proof}
%
%The above Lemma \ref{lem:noHall} immediately implies that the following statement is vacuously true.
%
%\begin{theorem}\label{thm:alternating}
%Let $n\in\ZZ_{\ge 5}$ and $\pi$ be a set of odd primes such that
%$\Al_n$ has a Hall $\pi$-subgroup. Then there exists
%$1_{\Al_n}\neq\chi\in\Irr_{\pi'}(\Al_n)$ with
%$\QQ(\chi)\subseteq\QQ(\zeta_p)$ for some $p\in\pi$.
%\end{theorem}
%
%Notice that the $(2\notin \pi)$-hypothesis in the statement of Theorem \ref{thm:alternating} can not be removed.
%In fact, as we already pointed out in the introduction, $\Al_7$ posses a Hall $\{2,3,5\}$-subgroup but the only $\{2,3,5\}'$-degree irreducible character of $\Al_7$ is the trivial character.
%Actually, this observation can be extended to any prime larger than $7$.
%Let $p\geq 11$ be a fixed prime number, and let $\pi$ be the set consisting of all prime numbers strictly smaller than $p$.
%In this case $\Al_{p-1}$ is a Hall $\pi$-subgroup of $\Al_p$ but $\mathrm{Irr}_{\pi'}(\Al_p)=\{1_{\Al_p}\}$. In fact any non-trivial irreducible character $\chi$ of $\pi'$-degree of $\Al_p$ would satisfy $\chi(1)=p$. This is not possible because for any $\zeta\in\mathrm{Irr}(\Sy_p)$ we have that either $\zeta(1)\leq p-1$ or that $\zeta(1)\geq \frac{p(p-3)}{2}$, by \cite{Rasala}. Since $p-3> 4$ we have that for any $\theta\in\mathrm{Irr}(\Al_p)$ we have that either $\theta(1)\leq p-1$ or that $\theta(1)>p$.


%%%%%%%%%%%%%%%%%%%%%%%
%%%%%%%%%%%%%%%%%%%%%%%
\section{Non-Abelian Simple Groups}

This section is devoted to proving Theorem~\ref{thm:main} for all
non-abelian simple groups. We first fix some notation. For a
positive integer $k$, we %%
%% Correction 36 [ ]
%% Annotated text: "ζk :“ e2πi{k<Caret> and we"
%% Comment: "" 
%% 
%% START of correction 36
let $\zeta_k:=e^{2\pi i/k}$ and %%
%% END of correction 36
we write $\QQ(\zeta_k)$ for the $k$th
cyclotomic field. As usual, $\Irr(G)$ denotes
the set of all irreducible characters of a group $G$, and
$\Irr_{\pi'}(G)$ denotes the subset consisting of those characters
whose degrees are not divisible by any prime in $\pi$. For
$\chi\in\Irr(G)$, we write $\QQ(\chi)$ for the field of values of
$\chi$, that is, the smallest extension of $\QQ$ containing all
values of $\chi$. Finally, for integers $x\le y$, we let
$[x,y]:=\{\,z\in\ZZ \mid x\le z\le y\,\}$.

As mentioned above, we aim at proving the following.

\begin{theorem}\label{thm:simple}
Let $S$ be a non-abelian simple group and $\pi$ be a set of odd
%%
%% Correction 32 [ ]
%% Annotated text: "Qpχq Ď Qpζpq<Caret> for some"
%% Comment: "" 
%% 
%% Correction 35 [ ]
%% Annotated text: "χ P Irrπ1pSq<Caret> such that"
%% Comment: "" 
%% 
%% START of corrections 32, 35
primes such that $S$ has a Hall $\pi$-subgroup. Then there exists
$1_S\neq\chi\in\Irr_{\pi'}(S)$ such that
$\QQ(\chi)\subseteq\QQ(\zeta_p)$ for %%
%% END of corrections 32, 35
some $p\in\pi$.
\end{theorem}


Hall $\pi$-subgroups for a set $\pi$ of odd primes are relatively
common in simple groups of Lie type, but are much more restricted in
alternating and sporadic groups. We begin by treating the
alternating and sporadic cases.

\begin{lemma}\label{lem:noHall}
%%
%% Correction 34 [ ]
%% Annotated text: "n P Ně5<Caret> and let"
%% Comment: "" 
%% 
%% START of correction 34
Let $n\in\mathbb{N}_{\geq 5}$ and %%
%% END of correction 34
let $\pi\subseteq [1,n]$ be a set
of prime numbers. Assume that $2\notin\pi$. Then the alternating
group $\Al_n$ admits a Hall $\pi$-subgroup if and only if $|\pi|=1$.
\end{lemma}

\begin{proof}
%%
%% Correction 33 [ ]
%% Annotated text: "|π| “ 1<Caret> the statement"
%% Comment: "" 
%% 
%% START of correction 33
If $|\pi|=1$ the %%
%% END of correction 33
statement is obviously implied by the Sylow theory.
For the other implication, assume that $\Al_n$ admits a Hall
$\pi$-subgroup $H$ and suppose for a contradiction that $|\pi|\geq
2$. %%
%% Correction 37 [ ]
%% Annotated text: ". , ptu<Caret> for some"
%% Comment: "" 
%% 
%% Correction 38 [ ]
%% Annotated text: "In particular let<Caret> π “"
%% Comment: "" 
%% 
%% Correction 39 [ ]
%% Annotated text: "2. In particular<Caret> let π"
%% Comment: "" 
%% 
%% START of corrections 37, 38, 39
In particular let $\pi=\{p_1,p_2,\ldots, p_t\}$ for %%
%% END of corrections 37, 38, 39
some odd
primes $p_1<p_2<\cdots<p_t$. %We proceed by distinguishing two cases,
%depending on the solvability of $H$.
Notice that $H$ is solvable, by the Feit--Thompson theorem
\cite{Thompson63}. By \cite{Hall1} we then know that $H$ admits %%
%% Correction 40 [ ]
%% Annotated text: "a tp1, p2u-<Caret>Hall subgroup K."
%% Comment: "" 
%% 
%% Correction 41 [ ]
%% Annotated text: "a tp1, p2u<Caret>-Hall subgroup K."
%% Comment: "" 
%% 
%% Correction 42 [ ]
%% Annotated text: "a tp1, p2u-Hall<Caret> subgroup K."
%% Comment: "" 
%% 
%% START of corrections 40, 41, 42
a
$\{p_1,p_2\}$-Hall subgroup %%
%% END of corrections 40, 41, 42
$K$. It follows that $K$ is a solvable
$\{p_1,p_2\}$-Hall subgroup of $\Al_n$. Moreover, since $2\notin
\{p_1,p_2\}$ we also have that $K$ is a solvable $\{p_1,p_2\}$-Hall
subgroup of the symmetric group $\Sy_n$. Using \cite[Theorem
A4]{Hall2}, we deduce that $|K|$ is even, and %%
%% Correction 43 [ ]
%% Annotated text: "P tp1, p2u<Caret> π. This"
%% Comment: "" 
%% 
%% Correction 44 [ ]
%% Annotated text: "and therefore 2<Caret> tp1, p2u"
%% Comment: "" 
%% 
%% START of corrections 43, 44
therefore $2\in
\{p_1,p_2\}\subseteq \pi$. %%
%% END of corrections 43, 44
This clearly contradicts our hypothesis.
%
%
%%\noindent\textit{Case} (i).
%(i) First assume that $H$ is a solvable group. By \cite{Hall1} we
%know that $H$ admits a $\{p_1,p_2\}$-Hall subgroup $K$. It follows
%that $K$ is a solvable $\{p_1,p_2\}$-Hall subgroup of $\Al_n$.
%Moreover, since $2\notin \{p_1,p_2\}$ we also have that $K$ is a
%solvable $\{p_1,p_2\}$-Hall subgroup of the symmetric group $\Sy_n$. Using \cite[Theorem
%A4]{Hall2}, we deduce that $|K|$ is even, and therefore $2\in
%\{p_1,p_2\}\subseteq \pi$. This clearly contradicts our hypothesis.
\end{proof}




We note that the use of the Feit--Thompson theorem can be avoided by
instead relying on the known classification of non-solvable Hall
subgroups of symmetric groups. Suppose that $H$ is a non-solvable
Hall $\pi$-subgroup of $\Al_n$. Since $2 \notin \pi$, it follows
that $H$ is also a %%
%% Correction 45 [ ]
%% Annotated text: "a non-solvable Hall<Caret> π-subgroup of"
%% Comment: "" 
%% 
%% START of correction 45
non-solvable Hall $\pi$%%
%% END of correction 45
-subgroup of $\Sy_n$. This
family of subgroups of symmetric groups is completely described in
\cite{Thompson}. In particular, either $H = \Sy_n$ or $n$ is a prime
number and $H = \Sy_{n-1}$. In both cases, we would have $2 \in
\pi$, which again contradicts our hypothesis.


%The above Lemma \ref{lem:noHall} immediately implies that the following statement is vacuously true.

\begin{proposition}\label{prop:alternatingandsporadic}
Theorem~\ref{thm:simple} holds when $S$ is an alternating group, a
sporadic group, or the %%
%% Correction 46 [ ]
%% Annotated text: "Tits group 2F4p2q1.<Caret>"
%% Comment: "" 
%% 
%% Correction 47 [ ]
%% Annotated text: "the Tits group<Caret> 2F4p2q1."
%% Comment: "" 
%% 
%% START of corrections 46, 47
Tits group ${}^2F_4(2)'$.
%Let $n\in\ZZ_{\ge 5}$ and $\pi$ be a set of odd primes such that
%$\Al_n$ has a Hall $\pi$-subgroup. Then there exists
%$1_{\Al_n}\neq\chi\in\Irr_{\pi'}(\Al_n)$ with
%$\QQ(\chi)\subseteq\QQ(\zeta_p)$ for some $p\in\pi$.
\end{proposition}

\begin{proof}
The %%
%% END of corrections 46, 47
case of alternating groups follows from Lemma~\ref{lem:noHall}
and Theorem~\ref{thm:Navarro-Tiep}. We now consider the sporadic
groups and the Tits group. By \cite[Theorem~6.14]{G86}, if $S$ has a
Hall $\pi$-subgroup for a set $\pi$ of odd primes, %%
%% Correction 69 [ ]
%% Annotated text: "then |π| ď<Caret> The result"
%% Comment: "" 
%% 
%% START of correction 69
then $|\pi| \le
2$. %%
%% END of correction 69
The result then follows from \cite[Theorem~2.1]{GHSV21}, except
possibly in the cases $(S,\pi)=({}^2F_4(2)',\{3,5\})$, $(J_4,
\{23,43\})$, or $(J_4,\{29,43\})$. However, a direct check using
\cite{Atl1} shows that %%
%% Correction 65 [ ]
%% Annotated text: "that in □<Caret> of these cases, S does not have a Hall π-subgroup."
%% Comment: "" 
%% 
%% Correction 66 [ ]
%% Annotated text: "that in each<Caret> of these cases, S does not have a Hall π-subgroup."
%% Comment: "" 
%% 
%% START of corrections 65, 66
in each of these cases, %%
%% END of corrections 65, 66
$S$ does not have a
Hall %%
%% Correction 67 [ ]
%% Annotated text: "a Hall π-subgroup<Caret>. □"
%% Comment: "" 
%% 
%% Correction 68 [ ]
%% Annotated text: "a Hall π-subgroup.<Caret> □"
%% Comment: "" 
%% 
%% START of corrections 67, 68
$\pi$-subgroup.
\end{proof}

%In particular we have that Theorem \ref{thm:simple} holds whenever $S$ is a simple alternating group.

\begin{remark}
The %%
%% END of corrections 67, 68
hypothesis $2\notin \pi$ in
Proposition~\ref{prop:alternatingandsporadic} can not be removed.
Indeed, as already pointed out in the introduction, $\Al_7$
possesses a Hall $\{2,3,5\}$-subgroup %%
%% Correction 55 [ ]
%% Annotated text: "5u-subgroup but t<Caret>he only t2,"
%% Comment: "" 
%% 
%% START of correction 55
but the only %%
%% END of correction 55
$\{2,3,5\}'$-degree irreducible character of $\Al_7$ is the trivial
%%
%% Correction 53 [ ]
%% Annotated text: "trivial character. This<Caret> observation can"
%% Comment: "" 
%% 
%% Correction 56 [ ]
%% Annotated text: "character. This o<Caret>bservation can be"
%% Comment: "" 
%% 
%% START of corrections 53, 56
character. This observation can %%
%% END of corrections 53, 56
be extended to any prime larger than
$7$. %%
%% Correction 54 [ ]
%% Annotated text: "p ě 11<Caret> be a"
%% Comment: "" 
%% 
%% START of correction 54
Let $p\geq 11$ be %%
%% END of correction 54
a fixed %%
%% Correction 57 [ ]
%% Annotated text: "fixed prime n<Caret>umber, and let"
%% Comment: "" 
%% 
%% START of correction 57
prime number, and %%
%% END of correction 57
let $\pi$ be the
set consisting of all prime numbers strictly %%
%% Correction 52 [ ]
%% Annotated text: "strictly smaller than<Caret> p. In"
%% Comment: "" 
%% 
%% Correction 58 [ ]
%% Annotated text: "smaller than p<Caret>. In this"
%% Comment: "" 
%% 
%% START of corrections 52, 58
smaller than $p$. In %%
%% END of corrections 52, 58
this case $\Al_{p-1}$ is a Hall $\pi$-subgroup of $\Al_p$ but
$\mathrm{Irr}_{\pi'}(\Al_p)=\{1_{\Al_p}\}$. In fact any non-trivial
irreducible character $\chi$ of $\pi'$-degree of $\Al_p$ would
satisfy $\chi(1)=p$. %%
%% Correction 51 [ ]
%% Annotated text: "p. This is<Caret> not possible"
%% Comment: "" 
%% 
%% START of correction 51
This is not %%
%% END of correction 51
possible because for any
$\zeta\in\mathrm{Irr}(\Sy_p)$ we have either $\zeta(1)\leq p-1$ %%
%% Correction 49 [ ]
%% Annotated text: "ppp´3q 2 ,<Caret> by [Ras77]."
%% Comment: "" 
%% 
%% Correction 50 [ ]
%% Annotated text: "ζp1q ě ppp´3q 2<Caret>, 2"
%% Comment: "" 
%% 
%% START of corrections 49, 50
or
$\zeta(1)\geq \frac{p(p-3)}{2}$, by %%
%% END of corrections 49, 50
\cite{Rasala}. Since $p-3> 4$ we
have that for any $\theta\in\mathrm{Irr}(\Al_p)$, either
$\theta(1)\leq p-1$ or $\theta(1)>p$%%
%% Correction 59 [ ]
%% Annotated text: "W<Caret>e now turn"
%% Comment: "" 
%% 
%% START of correction 59
.
\end{remark}

\smallskip

We now %%
%% END of correction 59
turn to the proof of Theorem~\ref{thm:simple} for simple
groups of Lie type $S \neq {}^2F_4(2)'$. By this, we mean simple
groups of the form $S = G/\bZ(G)$%%
%% Correction 60 [ ]
%% Annotated text: "“ G{ZpGq, where<Caret> G :“"
%% Comment: "" 
%% 
%% Correction 61 [ ]
%% Annotated text: "G :“ GF<Caret> is the"
%% Comment: "" 
%% 
%% Correction 62 [ ]
%% Annotated text: "GF is the<Caret> group of"
%% Comment: "" 
%% 
%% START of corrections 60, 61, 62
, where $G := \bG^F$ is the group %%
%% END of corrections 60, 61, 62
of fixed points of a simple, simply connected algebraic group $\bG$
%%
%% Correction 63 [ ]
%% Annotated text: "G defined over<Caret> an algebraically"
%% Comment: "" 
%% 
%% START of correction 63
defined over an %%
%% END of correction 63
algebraically closed field of characteristic $\ell$,
under a Steinberg %%
%% Correction 64 [ ]
%% Annotated text: "Steinberg endomorphism F<Caret> of G."
%% Comment: "" 
%% 
%% START of correction 64
endomorphism $F$ of %%
%% END of correction 64
$\bG$.
%; see
%\cite[Theorem~24.17]{malletesterman}.

The %%
%% Correction 48 [ ]
%% Annotated text: "case ℓR π<Caret> is easy."
%% Comment: "" 
%% 
%% START of correction 48
case $\ell\notin\pi$ is %%
%% END of correction 48
easy.

\begin{proposition}\label{prop:ell-notin-pi}
Theorem \ref{thm:simple} holds when $S$ is a simple group of Lie
type in characteristic $\ell$ and $\ell\notin\pi$.
\end{proposition}

\begin{proof}
The Steinberg character $\St_G$ of $G$ is trivial on $\bZ(G)$ and
thus can be viewed as a character of $S$. Moreover, $\St_G$ is
rational-valued and has degree a power of $\ell$; see
\cite[Proposition~3.4.10]{GeckMalle}.
\end{proof}

We therefore focus on the case $\ell\in\pi$. Let $\bB$ be an
$F$-stable Borel subgroup of $\bG$ and $\bT$ a (maximally-split)
$F$-stable maximal torus of $\bG$ inside $\bB$. Let $\Phi$ be the
root system of $\bG$ with respect to $\bT$ and $\bB$ and $\Phi^+$
the set of corresponding positive roots. Let $\bU$ be the product of
the root subgroups corresponding to the roots in $\Phi^+$. This
$\bU$ is indeed the unipotent radical of $\bB$. We have
$U:=\bU^F\in\Syl_\ell(G)$ and $B=UT$, %%
%% Correction 76 [ ]
%% Annotated text: "T :“ TF;<Caret> see [Car85,"
%% Comment: "" 
%% 
%% Correction 77 [ ]
%% Annotated text: "T :“ TF<Caret>; see [Car85,"
%% Comment: "" 
%% 
%% Correction 78 [ ]
%% Annotated text: "U :“ UF<Caret> and T"
%% Comment: "" 
%% 
%% START of corrections 76, 77, 78
where $U:=\bU^F$ and
$T:=\bT^F$; see %%
%% END of corrections 76, 77, 78
\cite[\S1.9 and \S2.9]{Carter85}.

Let $(\bG^*,F^*)$ be the pair dual to $(\bG,F)$, and %%
%% Correction 74 [ ]
%% Annotated text: ":“ G˚F ˚.<Caret> Note that"
%% Comment: "" 
%% 
%% Correction 75 [ ]
%% Annotated text: ":“ G˚F ˚<Caret>. Note that"
%% Comment: "" 
%% 
%% START of corrections 74, 75
set
$G^*:={\bG^*}^{F^*}$. Note %%
%% END of corrections 74, 75
that $S=[G^*,G^*]$, as $S$ is simple. Let
$\bT^*$ be an $F^*$-stable maximal torus of $\bG^*$ that is dual to
$\bT$ in the sense of \cite[Proposition~4.3.1]{Carter85}. Let
$\bB^*$ be an %%
%% Correction 73 [ ]
%% Annotated text: "an F ˚-<Caret>stable Borel subgroup"
%% Comment: "" 
%% 
%% START of correction 73
$F^*$-stable Borel %%
%% END of correction 73
subgroup of $\bG^*$ %%
%% Correction 71 [ ]
%% Annotated text: "G˚ containing T˚<Caret>. Write T"
%% Comment: "" 
%% 
%% START of correction 71
containing
$\bT^*$. Write %%
%% END of correction 71
$T^*:={\bT^*}^{F^*}$ %%
%% Correction 72 [ ]
%% Annotated text: ":“ B˚F ˚<Caret>. If U˚"
%% Comment: "" 
%% 
%% START of correction 72
and $B^*:={\bB^*}^{F^*}$. If %%
%% END of correction 72
$\bU^*$ denotes the unipotent radical %%
%% Correction 70 [ ]
%% Annotated text: "radical of B˚<Caret> and U"
%% Comment: "" 
%% 
%% START of correction 70
of $\bB^*$ and %%
%% END of correction 70
$U^*:={\bU^*}^{F^*}$, then $B^*=U^*T^*$.


\begin{lemma}\label{lem:Gross}
Let $S$ be a simple group of Lie type in characteristic $\ell$.
Suppose that $\pi$ is a set of odd primes containing $\ell$ (in
particular, $\ell$ is odd) and that $S$ has a Hall $\pi$-subgroup.
Then $B^*$ contains a Hall $\pi$-subgroup of $G^*$.
\end{lemma}

\begin{proof}
Note that $\bZ(G)\le T\le B$. By \cite[Theorem~3.2]{G86}, the
quotient $B/\bZ(G)$, which may be regarded as a Borel subgroup of
$S$, contains a Hall $\pi$-subgroup of $S$. Equivalently, $|G:B| =
|S : B/\bZ(G)|$ is a $\pi'$-number. By Corollary~4.4.2 and
Proposition~4.4.4 of \cite{Carter85}, we have $|G^*|=|G|$ and
$|T^*|=|T|$, and it follows that $|U^*|=|U|$ and $|B|=|B^*|$. (Note
that $U$ and $U^*$ have orders equal to the $\ell$-parts of $|G|$
and $|G^*|$, respectively.) Now we have that $|G^*:B^*|$ is a
$\pi'$-number. As $B^*$ is solvable, it therefore contains a Hall
$\pi$-subgroup of $G^*$.
\end{proof}

\begin{lemma}\label{lem:Gross2}
Let $S\neq {}^2F_4(2)'$ be a simple group of Lie type in
characteristic $\ell$. To prove Theorem~\ref{thm:simple} for $S$, it
is sufficient to assume that $\pi\cap\pi(T^*)\neq\emptyset$.
\end{lemma}

\begin{proof}
By Proposition~\ref{prop:ell-notin-pi}, we may assume that
$\ell\in\pi$. If $\pi\cap \pi(S)=\{\ell\}$, then
Theorem~\ref{thm:simple} follows from
Theorem~\ref{thm:Navarro-Tiep}. We therefore may assume that
$|\pi\cap \pi(S)|\geq 2$. Since $S$ has a Hall $\pi$-subgroup, by
Lemma~\ref{lem:Gross}, the Borel subgroup $B^*$ contains a Hall
$\pi$-subgroup of $G^*$. In particular, $|\pi\cap \pi(B^*)|\geq 2$,
and hence there exists at least one prime lying in both $\pi$ and
$\pi(T^*)$.
\end{proof}



Recall that $S=G/\bZ(G)$, so the irreducible characters of $S$ are
precisely those irreducible characters of $G$ whose kernel contains
$\bZ(G)$. The set $\Irr(G)$ admits a natural partition into Lusztig
series $\EC(G,s)$ indexed by $G^\ast$-conjugacy classes of
semisimple elements $s\in G^\ast$. The series $\EC(G,s)$ consists of
those irreducible characters of $G$ that are constituents of some
Deligne--Lusztig character $R_{\bS}^{\bG}(\theta)$, where $\bS$ is
an $F$-stable maximal torus of $\bG$ and $\theta\in\Irr(\bS^F)$ is
such that the geometric conjugacy class of the pair $(\bS,\theta)$
corresponds to the $G^\ast$-conjugacy class containing $s$. See
\cite[\S2.6]{GeckMalle}.

We will look for the desired character among the so-called
\emph{semisimple characters}. Recall that \(\rho \in \Irr(G)\) is
called a semisimple character if the average value of \(\rho\) on
\(\mathfrak{C}^F\) is nonzero, where \(\mathfrak{C}\) is the
conjugacy class of \(\bG\) consisting of regular unipotent elements.
For each semisimple element \(s \in G^*\), the Lusztig series
\(\EC(G,s)\) contains one or more semisimple characters, all of
which have %%
%% Correction 79 [ ]
%% Annotated text: "|G˚ : CG˚psq|ℓ1.<Caret>"
%% Comment: "" 
%% 
%% START of correction 79
degree
\[
|G^* : \bC_{G^*}(s)|_{\ell'}.
\]
In %%
%% END of correction 79
fact, when \(\bC_{\bG^*}(s)\) is connected, \(\EC(G,s)\) contains
exactly one semisimple character; see \cite[p.~171]{GeckMalle}. We
record a well-known fact about these characters below. Here we use
$\ord(g)$ to denote the order of a group element $g$.

\begin{lemma}\label{lem:1}
Assume the above notation. Assume furthermore that $\bC_{\bG^*}(s)$
is connected. Then the unique semisimple character, say $\chi_s$, in
the series \(\EC(G,s)\) is trivial on $\bZ(G)$ if and only if $s\in
S$. In such situation, $\QQ(\chi_s)\subseteq \QQ(\zeta_{\ord(s)})$.
\end{lemma}

\begin{proof}
The first part follows from, for instance,
\cite[Proposition~24.21]{malletesterman}, \cite[Lemma~2.2]{M07}, and
\cite[Lemma~5.8]{Hung22}. The latter part is
\cite[Lemma~4.2]{GHSV21}.
\end{proof}

We handle the linear and unitary groups separately. To unify
notation, we write $S=\PSL_n^\epsilon(q)$, where the superscript
$\epsilon=+1$ corresponds to the linear groups and $\epsilon=-1$ to
the unitary groups. We use analogous notation for the related
groups, for example $G=\SL_n^\epsilon(q)$ and
$G^*=\PGL_n^\epsilon(q)$. It is more convenient to first study the
characters of $\widetilde{G}:=\GL_n^\epsilon(q)$ and then analyze
those of $G=\SL_n^\epsilon(q)$ as irreducible constituents of
restricted characters. Note that $\widetilde{G}$ is self-dual, and
we will identify $\widetilde{G}$ with its dual group. Let
$\pi:\widetilde{G}\to G^*$ be the natural projection from
$\widetilde{G}$ to $G^*$.

Let $\widetilde{s}$ be a semisimple element of $\widetilde{G}$.
Since the ambient algebraic group of $\widetilde{G}$ has connected
center, the Lusztig series
$\mathcal{E}(\widetilde{G},\widetilde{s})$ contains a unique
semisimple character, which we denote by $\chi_{\widetilde{s}}$.

\begin{lemma}\label{lem:numberofconstituents}
The number of irreducible constituents of the restriction of
semisimple character $\chi_{\widetilde{s}}$ to $G$ divides
$\gcd(\ord(\widetilde{s}),n,q-\epsilon)$.
\end{lemma}

\begin{proof}
Set $s:=\pi(\widetilde{s})$. Recall that
$\chi_{\widetilde{s}}(1)=|\widetilde{G}:\bC_{\widetilde{G}}(\widetilde{s})|_{\ell'}$.
The restriction of $\chi_{\widetilde{s}}$ from $\widetilde{G}$ to
$G$ is multiplicity-free and its irreducible constituents are
precisely the semisimple characters of the Lusztig series
\(\EC(G,s)\) by \cite[Corollary~2.6.18]{GeckMalle}. Each of these
constituent has degree $|G^* : \bC_{G^*}(s)|_{\ell'}$. Therefore,
the number of irreducible constituents of the restriction %%
%% Correction 80 [ ]
%% Annotated text: ": π´1pCG˚psqq|ℓ1 “<Caret>|π : C"
%% Comment: "" 
%% 
%% Correction 81 [ ]
%% Annotated text: ": C rGpsq|ℓ1.<Caret>"
%% Comment: "" 
%% 
%% Correction 82 [ ]
%% Annotated text: "π´1pCG˚psqq|ℓ1 “ |π 1pCG˚psqq :<Caret>C : rGprsq|ℓ1."
%% Comment: "" 
%% 
%% Correction 83 [ ]
%% Annotated text: "|G˚ : CG˚psq|ℓ1<Caret> “ |"
%% Comment: "" 
%% 
%% START of corrections 80, 81, 82, 83
is
\[
\frac{|\widetilde{G}:\bC_{\widetilde{G}}(\widetilde{s})|_{\ell'}}{|G^*
:
\bC_{G^*}(s)|_{\ell'}}=\frac{|\widetilde{G}:\bC_{\widetilde{G}}(\widetilde{s})|_{\ell'}}{|\widetilde{G}
:
\pi^{-1}(\bC_{G^*}(s))|_{\ell'}}=|\pi^{-1}(\bC_{G^*}(s)):\bC_{\widetilde{G}}(\widetilde{s})|_{\ell'}.
\]




Let %%
%% END of corrections 80, 81, 82, 83
$Z$ be the (cyclic) subgroup of the multiplicative group
$\FF_{q^2-1}^\times$ of order $q-\epsilon$. Consider the mapping
%%
%% Correction 94 [ ]
%% Annotated text: "Ñ Z def<Caret>ned by grsg´1"
%% Comment: "" 
%% 
%% Correction 95 [ ]
%% Annotated text: "Ñ Z def<Caret>ned by grsg´1"
%% Comment: "" 
%% 
%% Correction 96 [ ]
%% Annotated text: "Ñ Z de<Caret>fined by grsg´1"
%% Comment: "" 
%% 
%% START of corrections 94, 95, 96
$\kappa: \pi^{-1}(\bC_{G^*}(s))\to Z$ defined by %%
%% END of corrections 94, 95, 96
$g\widetilde{s}g^{-1}=\kappa(g)\widetilde{s}$. This is indeed a
homomorphism with $\Ker(\kappa)=\bC_{\widetilde{G}}(\widetilde{s})$.
It follows that the index
$|\pi^{-1}(\bC_{G^*}(s)):\bC_{\widetilde{G}}(\widetilde{s})|$ is
equal to the order of the image $\Im(\kappa)$ of $\kappa$.
Therefore, the lemma follows once we show %%
%% Correction 84 [ ]
%% Annotated text: "q ´ ϵq.<Caret>"
%% Comment: "" 
%% 
%% START of correction 84
that
\[
|\Im(\kappa)| \text{ divides }
\gcd(\ord(\widetilde{s}),n,q-\epsilon).
\]


First, %%
%% END of correction 84
it %%
%% Correction 85 [ ]
%% Annotated text: "it is clear<Caret> that |Impκq|"
%% Comment: "" 
%% 
%% Correction 86 [ ]
%% Annotated text: "that |Impκq| |<Caret> pq ´"
%% Comment: "" 
%% 
%% Correction 87 [ ]
%% Annotated text: "pq ´ ϵq<Caret>. Now let g be an arbitrary element of π´1pCG˚psqq. We have detprsq “ “ detpκpgqrsq"
%% Comment: "" 
%% 
%% START of corrections 85, 86, 87
is clear that $|\Im(\kappa)|\mid (q-\epsilon)$. Now %%
%% END of corrections 85, 86, 87
let
$g$ be an arbitrary element of $\pi^{-1}(\bC_{G^*}(s))$. We have
$\det(\widetilde{s})=\det(g\widetilde{s}g^{-1})=\det(\kappa(g)\widetilde{s})=\kappa(g)^n\det(\widetilde{s})$,
and hence $\kappa(g)^n=1$; equivalently, $\ord(\kappa(g))\mid n$.
Finally, we observe that
$\ord(\widetilde{s})=\ord(g\widetilde{s}g^{-1})=\ord(\kappa(g)\widetilde{s})=\lcm(\ord(\kappa(g)),\widetilde{s})$,
which implies that $\ord(\kappa(g))\mid \ord(\widetilde{s})$.
%\[\{\delta^{(q-1)/p}, \delta^{(q-1)(p-1)/p},
%1^{\,n-2}\}=\{\kappa(g)\delta^{(q-1)/p},
%\kappa(g)\delta^{(q-1)(p-1)/p}, \kappa(g)^{\,n-2}\},\] where
%$\kappa(g)\in\bZ(\widetilde{G})$ is identified with the
%corresponding element in $\FF_q^\times$. This happens if and only if
%$\kappa(g)=1$. We have shown that $\pi^{-1}(\bC_{G^*}(s))=
%\bC_{\widetilde{G}}(\widetilde{s})$, and the claim follows.
%Now we have
%\[
%|G^*:\bC_{G^*}(s)|_{\ell'}=|G^*:\bC_{\widetilde{G}}(\widetilde{s})/\bZ(\widetilde{G})|_{\ell'}
%=|\widetilde{G}:\bC_{\widetilde{G}}(\widetilde{s})|_{\ell'},
%\]
%and therefore the semisimple characters in the series
%$\mathcal{E}(G,s)$ have the same degree as the (only) semisimple
%character $\chi_{\widetilde{s}}$ in
%$\mathcal{E}(\widetilde{G},\widetilde{s})$. However, by
%\cite[Corollary~2.6.18]{GeckMalle}, semisimple characters in
%$\mathcal{E}(G,s)$ are irreducible constituents of the restriction
%of $\chi_{\widetilde{s}}$ from $\widetilde{G}$ to $G$. We deduce
%that $\mathcal{E}(G,s)$ has a unique semisimple character, namely
%the restriction of $\chi_{\widetilde{s}}$ to $G$. Let us denote this
%restriction by $\psi$.
\end{proof}

The proof of Lemma~\ref{lem:numberofconstituents} also yeilds the
following well-known result; however, we only require the %%
%% Correction 88 [ ]
%% Annotated text: "the first statement.<Caret>"
%% Comment: "" 
%% 
%% Correction 89 [ ]
%% Annotated text: "the first statement<Caret>."
%% Comment: "" 
%% 
%% START of corrections 88, 89
first
statement.

\begin{lemma}\label{lem:2}
Let %%
%% END of corrections 88, 89
$s := \pi(\widetilde{s}) \in G^*$, where $\widetilde{s}$ is a
semisimple element of $\widetilde{G}$. Suppose that the multiset of
eigenvalues of $\widetilde{s}$ is not invariant under multiplication
by every nontrivial root of unity. Then the centralizer
$\bC_{\bG^*}(s)$ is connected. Moreover, if the multiset of
eigenvalues of $\widetilde{s}$ is not invariant under multiplication
by every nontrivial element of the subgroup of $\FF_{q^2-1}^\times$
of order $q-\epsilon$, then $\chi_{\widetilde{s}}$ restricts
irreducibly to $G$.
\end{lemma}

\begin{proof}
For every $g \in \pi^{-1}\bigl(\bC_{G^*}(s)\bigr)$, we have
$g\widetilde{s}g^{-1} = \kappa(g)\widetilde{s}$, and hence
$g\widetilde{s}g^{-1}$ and $\widetilde{s}$ have the same
eigenvalues. Therefore, if the multiset of eigenvalues of
$\widetilde{s}$ is not invariant under multiplication by every
nontrivial element of the subgroup of $\FF_{q^2-1}^\times$ of order
$q-\epsilon$, the homomorphism $\kappa$ must have trivial image,
which proves the second statement.

A similar argument, now in the setting of algebraic groups
$\widetilde{\bG} := \GL^\epsilon_n(\overline{\FF}_{\ell})$ and
${\bG^*} := \PGL^\epsilon_n(\overline{\FF}_{\ell})$ instead, shows
that if the multiset of eigenvalues of $\widetilde{s}$ is not
invariant under multiplication by every nontrivial root of unity,
then $\bC_{\bG^*}(s)$ is the image of the connected group
$\bC_{\widetilde{\bG}}(\widetilde{s})$ under $\pi$ and hence is
itself connected.
\end{proof}


\begin{proposition}\label{prop:linear-unitary}
Theorem~\ref{thm:simple} holds for the simple groups $S=\PSL_n(q)$
and $\PSU_n(q)$, where $n\geq 2$ and $q$ is a prime power.
\end{proposition}

\begin{proof}
Recall that $T^*$ is a maximal torus contained in a Borel subgroup
$B^*$ of $G^*$. By Lemma~\ref{lem:Gross2}, we may assume that $\pi$
and $\pi(T^*)$ share a common prime, say $p$. In the linear case, we
%%
%% Correction 90 [ ]
%% Annotated text: "pq ´ 1qn´1<Caret>,"
%% Comment: "" 
%% 
%% Correction 91 [ ]
%% Annotated text: "pq ´ 1qn´1,<Caret>"
%% Comment: "" 
%% 
%% START of corrections 90, 91
have
\[
|T^*| = (q-1)^{n-1},
\]
so %%
%% END of corrections 90, 91
$p$ divides $q-1$. In the unitary %%
%% Correction 92 [ ]
%% Annotated text: "n is odd even<Caret>, n"
%% Comment: "" 
%% 
%% Correction 93 [ ]
%% Annotated text: "pq2 ´ 1qp q{ , 1qpn´2q{2p<Caret>q ´ if n is odd,"
%% Comment: "" 
%% 
%% START of corrections 92, 93
case,
\[
|T^*| =
\begin{cases}
(q^2-1)^{(n-1)/2}, & \text{if $n$ is odd},\\[2pt]
(q^2-1)^{(n-2)/2}(q-1), & \text{if $n$ is even},
\end{cases}
\]
and %%
%% END of corrections 92, 93
hence $p$ divides $q^2-1$ (see \cite[p.~74]{Carter85}). Since
$B^*$ contains a Hall $\pi$-subgroup of $G^*$ by
Lemma~\ref{lem:Gross}, a comparison of the orders of $G^*$ and $B^*$
shows that $p$ cannot divide $q+1$. Therefore, as in the linear
case, we again conclude that $p$ divides $q-1$.



We first consider the case $\epsilon=1$. Let $\delta$ be a generator
of $\FF_{q}^\times$, and define the semisimple %%
%% Correction 97 [ ]
%% Annotated text: "n´2˘ P G.<Caret>"
%% Comment: "" 
%% 
%% START of correction 97
element
\[
\widetilde{s} := \diag\!\bigl( \delta^{(q-1)/p},
\delta^{(q-1)(p-1)/p}, 1^{\,n-2} \bigr) \in \widetilde{G}.
\]
Let %%
%% END of correction 97
$s$ denote the image $\pi(\widetilde{s})$ of $\widetilde{s}$
%%
%% Correction 98 [ ]
%% Annotated text: "rG Ñ G˚<Caret>. Then"
%% Comment: "" 
%% 
%% Correction 99 [ ]
%% Annotated text: "Ñ G˚. Then<Caret>"
%% Comment: "" 
%% 
%% Correction 100 [ ]
%% Annotated text: "ordprsq “ p.<Caret>"
%% Comment: "" 
%% 
%% START of corrections 98, 99, 100
under $\pi:\widetilde{G}\to G^*$. Then
\[
\ord(s) = \ord(\widetilde{s}) = p.
\]
Taking %%
%% END of corrections 98, 99, 100
conjugation if necessary, we may assume that $s\in T^*$.
Moreover, as $\det(\widetilde{s})=\delta^{q-1}=1$, we %%
%% Correction 101 [ ]
%% Annotated text: "rs P G<Caret> and s"
%% Comment: "" 
%% 
%% Correction 102 [ ]
%% Annotated text: "X T ˚.<Caret> Note that"
%% Comment: "" 
%% 
%% Correction 103 [ ]
%% Annotated text: "X T ˚<Caret>. Note that"
%% Comment: "" 
%% 
%% START of corrections 101, 102, 103
have
\[
\widetilde{s}\in G \quad\text{and}\quad s\in S\cap T^*.
\]

Note %%
%% END of corrections 101, 102, 103
that the case $(n,p)=(3,3)$ cannot occur. (Otherwise, by
Lemma~\ref{lem:Gross}, the Borel subgroup $B^*$, of order
$q^3(q-1)^2$, would contain a Sylow $3$-subgroup of $G^*$ whose
order $q^3(q^2-1)(q^3-1)$ has larger $3$-part, leading to a
contradiction.) It is then easy to check that the multiset of
eigenvalues of $\widetilde{s}$ is not invariant under multiplication
by every nontrivial root of unity. By Lemma~\ref{lem:2}, %%
%% Correction 141 [ ]
%% Annotated text: "2.10, it fol<Caret>lows that CG˚psq"
%% Comment: "" 
%% 
%% Correction 142 [ ]
%% Annotated text: "2.10, it fo<Caret>llows that CG˚psq"
%% Comment: "" 
%% 
%% Correction 143 [ ]
%% Annotated text: "2.10, it foll<Caret>ows that CG˚psq"
%% Comment: "" 
%% 
%% START of corrections 141, 142, 143
it follows
that %%
%% END of corrections 141, 142, 143
$\bC_{\bG^*}(s)$ is connected, and we let $\chi_s$ be the
unique semisimple character in the series $\EC(G,s)$. Using
Lemma~\ref{lem:1}, we deduce %%
%% Correction 104 [ ]
%% Annotated text: "Qpζordpsqq “ Qpζpq.<Caret>"
%% Comment: "" 
%% 
%% START of correction 104
that
\[\QQ(\chi_{s})\subseteq\QQ(\zeta_{\ord(s)})=\QQ(\zeta_p).\] Moreover, %%
%% END of correction 104
since $s\in T^*$ and $T^*$ is abelian, we %%
%% Correction 105 [ ]
%% Annotated text: "˚ ď CG˚psq.<Caret>"
%% Comment: "" 
%% 
%% Correction 106 [ ]
%% Annotated text: ": T ˚|,<Caret> and hence"
%% Comment: "" 
%% 
%% Correction 107 [ ]
%% Annotated text: "It follows that<Caret> |G˚ :"
%% Comment: "" 
%% 
%% Correction 108 [ ]
%% Annotated text: "It follow<Caret>s that |G˚"
%% Comment: "" 
%% 
%% Correction 109 [ ]
%% Annotated text: "It follows<Caret> that |G˚"
%% Comment: "" 
%% 
%% Correction 110 [ ]
%% Annotated text: "It fo<Caret>llows that |G˚"
%% Comment: "" 
%% 
%% Correction 111 [ ]
%% Annotated text: "It fol<Caret>lows that |G˚"
%% Comment: "" 
%% 
%% Correction 112 [ ]
%% Annotated text: "˚|, and h<Caret>ence χsp1q “"
%% Comment: "" 
%% 
%% Correction 113 [ ]
%% Annotated text: "˚|, and he<Caret>nce χsp1q “"
%% Comment: "" 
%% 
%% Correction 114 [ ]
%% Annotated text: "˚|, and hen<Caret>ce χsp1q “"
%% Comment: "" 
%% 
%% Correction 115 [ ]
%% Annotated text: "˚|, and henc<Caret>e χsp1q “"
%% Comment: "" 
%% 
%% Correction 116 [ ]
%% Annotated text: "˚|, and hence<Caret> χsp1q “"
%% Comment: "" 
%% 
%% Correction 117 [ ]
%% Annotated text: "T ˚|, and<Caret> hence χsp1q"
%% Comment: "" 
%% 
%% Correction 118 [ ]
%% Annotated text: "T ˚|, an<Caret>d hence χsp1q"
%% Comment: "" 
%% 
%% Correction 119 [ ]
%% Annotated text: "T ˚|, a<Caret>nd hence χsp1q"
%% Comment: "" 
%% 
%% START of corrections 105, 106, 107, 108, 109, 110, 111, 112, 113, 114, 115, 116, 117, 118, 119
have \[T^*\le
\bC_{\bG^*}(s).\] It follows that \[|G^*:\bC_{\bG^*}(s)| \text{
divides } |G^*:T^*|,\] and hence
\[
\chi_s(1)=|G^*:\bC_{G^*}(s)|_{\ell'} \text{ divides }
|G^*:T^*|_{\ell'}.
\]
As %%
%% END of corrections 105, 106, 107, 108, 109, 110, 111, 112, 113, 114, 115, 116, 117, 118, 119
$|G^*:T^*|_{\ell'}=|G^*:B^*|$ is a $\pi'$-number %%
%% Correction 140 [ ]
%% Annotated text: "π1-number by Lem<Caret>ma 2.6, we"
%% Comment: "" 
%% 
%% START of correction 140
by
Lemma~\ref{lem:Gross}, we %%
%% END of correction 140
conclude that $\chi_s(1)$ is also a
$\pi'$-number. Also, from the fact that $s\in S=[G^*,G^*]$, we know
that $\chi_{s}$ is trivial %%
%% Correction 120 [ ]
%% Annotated text: "trivial on Z<Caret>pGq. This completes"
%% Comment: "" 
%% 
%% START of correction 120
on $\bZ({G})$. %%
%% END of correction 120
This completes the proof
for linear %%
%% Correction 121 [ ]
%% Annotated text: "linear groups. I<Caret>t remains to"
%% Comment: "" 
%% 
%% Correction 122 [ ]
%% Annotated text: "groups. It r<Caret>emains to consider"
%% Comment: "" 
%% 
%% Correction 123 [ ]
%% Annotated text: "groups. It rem<Caret>ains to consider"
%% Comment: "" 
%% 
%% Correction 124 [ ]
%% Annotated text: "groups. It re<Caret>mains to consider"
%% Comment: "" 
%% 
%% Correction 125 [ ]
%% Annotated text: "groups. It rema<Caret>ins to consider"
%% Comment: "" 
%% 
%% Correction 126 [ ]
%% Annotated text: "groups. It remain<Caret>s to consider"
%% Comment: "" 
%% 
%% Correction 127 [ ]
%% Annotated text: "It remains t<Caret>o consider the"
%% Comment: "" 
%% 
%% Correction 128 [ ]
%% Annotated text: "remains to co<Caret>nsider the case"
%% Comment: "" 
%% 
%% Correction 129 [ ]
%% Annotated text: "remains to c<Caret>onsider the case"
%% Comment: "" 
%% 
%% Correction 130 [ ]
%% Annotated text: "remains to con<Caret>sider the case"
%% Comment: "" 
%% 
%% Correction 131 [ ]
%% Annotated text: "remains to consid<Caret>er the case"
%% Comment: "" 
%% 
%% Correction 132 [ ]
%% Annotated text: "remains to consid<Caret>er the case"
%% Comment: "" 
%% 
%% Correction 133 [ ]
%% Annotated text: "remains to consi<Caret>der the case"
%% Comment: "" 
%% 
%% Correction 134 [ ]
%% Annotated text: "remains to cons<Caret>ider the case"
%% Comment: "" 
%% 
%% Correction 135 [ ]
%% Annotated text: "remains to conside<Caret>r the case"
%% Comment: "" 
%% 
%% Correction 136 [ ]
%% Annotated text: "to consider t<Caret>he case ϵ"
%% Comment: "" 
%% 
%% Correction 137 [ ]
%% Annotated text: "to consider the<Caret> case ϵ"
%% Comment: "" 
%% 
%% Correction 138 [ ]
%% Annotated text: "to consider th<Caret>e case ϵ"
%% Comment: "" 
%% 
%% Correction 139 [ ]
%% Annotated text: "consider the c<Caret>ase ϵ “"
%% Comment: "" 
%% 
%% START of corrections 121, 122, 123, 124, 125, 126, 127, 128, 129, 130, 131, 132, 133, 134, 135, 136, 137, 138, 139
groups.


It remains to consider the case $\epsilon=-1$ %%
%% END of corrections 121, 122, 123, 124, 125, 126, 127, 128, 129, 130, 131, 132, 133, 134, 135, 136, 137, 138, 139
and $p \mid (q-1)$.
Let \(\widetilde{t}\) be a semisimple element of order \(p\) in a
maximal torus of order $q^2-1$ of \(\GU_2(q)\). Since \(p \mid
(q-1)\) and \(p\) is odd, we have \(\gcd(p,q+1)=1\), and hence
\(\widetilde{t}\in \SU_2(q)\). Define the semisimple element
\[
\widetilde{s}:=\diag\left(\widetilde{t}, I_{n-2}\right)\in
\widetilde{G},
\]
and consider the corresponding semisimple character
\(\chi_{\widetilde{s}}\in \Irr(\widetilde{G})\). Note that
$\widetilde{s}$ belongs to a maximal torus of $\widetilde{G}$ of
order $(q^2-1)^{(n-1)/2}(q+1)$ if $n$ is odd, and of %%
%% Correction 144 [ ]
%% Annotated text: "of order pq2´1qn{2<Caret>"
%% Comment: "" 
%% 
%% Correction 145 [ ]
%% Annotated text: "n is even.<Caret>"
%% Comment: "" 
%% 
%% Correction 146 [ ]
%% Annotated text: "if n<Caret> is even."
%% Comment: "" 
%% 
%% START of corrections 144, 145, 146
order
$(q^2-1)^{n/2}$ if $n$ is even.

Note %%
%% END of corrections 144, 145, 146
also that \(\ord(\widetilde{s})=\ord(\widetilde{t})=p\), which
does not divide $q+1$. Lemma~\ref{lem:numberofconstituents} then
implies that \(\chi_{\widetilde{s}}\) restricts irreducibly to
\(G\). Arguments similar to those in the linear case show that this
restriction is trivial on \(\bZ(G)\), has \(\pi'\)-degree, and takes
values in \(\QQ(\zeta_p)\). This completes the proof.
\end{proof}


We can now finish the proof of Theorem~\ref{thm:simple}.

\begin{proof}[Proof of Theorem~\ref{thm:simple}]
By Propositions~\ref{prop:alternatingandsporadic} and
\ref{prop:linear-unitary}, and the classification of finite simple
groups, we may assume %%
%% Correction 160 [ ]
%% Annotated text: "S ‰ 2F4p2q1<Caret> is a"
%% Comment: "" 
%% 
%% START of correction 160
that $S \neq {}^2F_4(2)'$ is %%
%% END of correction 160
a simple group of
Lie type not of type $A$. By Lemma~\ref{lem:Gross2}, there exists at
least one prime lying in both $\pi$ and $\pi(T^*)$. As before, let
$p$ be such a prime.

Note that
\[
|T^*:(T^*\cap S)| = |T^*S:S| \text{ divides } |G^*:S|
\]
and $|G^*:S|$ is the order of the group of \emph{diagonal
automorphisms} of $S$. On the other hand, the order of $T^*$ is
given %%
%% Correction 147 [ ]
%% Annotated text: "1 ˘ ,"
%% Comment: "" 
%% 
%% START of correction 147
by
\[
|T^*|=\prod_{\mathfrak{o}\in \mathcal{O}}
\bigl(q^{|\mathfrak{o}|}-1\bigr),
\]
where %%
%% END of correction 147
$\mathcal{O}$ is the set of $F^*$-orbits on the simple roots
of the root system of $\bG^*$ associated with $\bB^*$ and $\bT^*$,
and $q$ is the absolute value of the eigenvalues of $F^*$ on the
character group of $\bT^*$; see \cite[p.~74]{Carter85}. A
straightforward case-by-case check, using \cite[Table~5]{Atl1} for
the size of the group of diagonal automorphisms and
\cite[\S1.19]{Carter85} for the sizes of the $F^*$-orbits, reveals
that if an odd prime $p$ is a divisor of $|G^*:S|$, then indeed
$(|G^*:S|_p)^2$ divides $|T^*|$, and it follows that $p$ also
divides $|T^*\cap S|$. (The only exception is the case $S=\PSU_3(q)$
with $(q+1)_3=3$, but this was already excluded using
Proposition~\ref{prop:linear-unitary}.) Therefore, we may and will
assume from now on that \[p\in \pi\cap \pi(T^*\cap S).\]
Consequently, there exists a semisimple %%
%% Correction 153 [ ]
%% Annotated text: "s P G˚<Caret> such that"
%% Comment: "" 
%% 
%% Correction 154 [ ]
%% Annotated text: "s P G<Caret>˚ such that"
%% Comment: "" 
%% 
%% START of corrections 153, 154
element $s\in G^*$ such %%
%% END of corrections 153, 154
%%
%% Correction 152 [ ]
%% Annotated text: "ordpsq “ p.<Caret>"
%% Comment: "" 
%% 
%% START of correction 152
that
\[
s\in T^*\cap S \quad\text{and}\quad \ord(s)=p.
\]

Suppose %%
%% END of correction 152
first that $p\nmid |\bZ(\bG)|$. Then $\bC_{\bG^*}(s)$ is
connected by \cite[Corollary~4.6]{Borel70}. Arguing as in the proof
of Proposition~\ref{prop:linear-unitary} and applying
Lemmas~\ref{lem:Gross} and \ref{lem:1}, we have that the unique
semisimple character $\chi_s$ in $\EC(G,s)$ has $\pi'$-degree, is
trivial on $\bZ(G)$, and satisfies $\QQ(\chi_s)\subseteq
\QQ(\zeta_p)$.

Next suppose that $p\mid |\bZ(\bG)|$. By
\cite[Theorem~1.12.5]{GLS98} and the assumption that $p$ is odd, we
are left with only the %%
%% Correction 151 [ ]
%% Annotated text: "p “ 3<Caret> and G"
%% Comment: "" 
%% 
%% START of correction 151
case $p=3$ and %%
%% END of correction 151
$\bG$ to be of type $E_6$.
%%
%% Correction 148 [ ]
%% Annotated text: "ϵ P t˘1u<Caret>, where ϵ"
%% Comment: "" 
%% 
%% Correction 149 [ ]
%% Annotated text: "ϵ P t˘1u,<Caret> where ϵ"
%% Comment: "" 
%% 
%% Correction 150 [ ]
%% Annotated text: "ϵ “ 1<Caret> corresponds to"
%% Comment: "" 
%% 
%% Correction 157 [ ]
%% Annotated text: "“ Eϵ Eϵ 6pqq<Caret> with ϵ"
%% Comment: "" 
%% 
%% START of corrections 148, 149, 150, 157
Here, $G=E_6^\epsilon(q)$ with $\epsilon\in\{\pm1\}$, where %%
%% END of corrections 148, 149, 150, 157
$\epsilon=1$ corresponds to the untwisted groups %%
%% Correction 158 [ ]
%% Annotated text: "ϵ “ ´1<Caret> to the"
%% Comment: "" 
%% 
%% Correction 159 [ ]
%% Annotated text: "ϵ “ ´<Caret>1 to the"
%% Comment: "" 
%% 
%% START of corrections 158, 159
and $\epsilon=-1$
to %%
%% END of corrections 158, 159
the twisted groups, and $|\bZ(G)|=\gcd(3,q-\epsilon)=3$.

As mentioned above, all the semisimple characters in the Lusztig
series $\EC(G,s)$ indexed by $s$ have the same degree
$|G^*:\bC_{G^*}(s)|_{\ell'}$. This is a $\pi'$-number due to the
fact that %%
%% Correction 155 [ ]
%% Annotated text: "P T ˚<Caret>, as argued"
%% Comment: "" 
%% 
%% Correction 156 [ ]
%% Annotated text: "P T ˚,<Caret> as argued"
%% Comment: "" 
%% 
%% START of corrections 155, 156
$s\in T^*$, as %%
%% END of corrections 155, 156
argued in the proof of
Proposition~\ref{prop:linear-unitary}. Moreover, as noted in the
proof of \cite[Proposition~4.5]{GHSV21}, these semisimple characters
take integer values on unipotent elements and therefore have field
of values contained in $\QQ(\zeta_3)$, by \cite[Lemma~4.3]{GHSV21}.
Finally, since their degrees are coprime to $3$, their restrictions
%%
%% Correction 180 [ ]
%% Annotated text: "restrictions to ZpGq<Caret> are multiples"
%% Comment: "" 
%% 
%% START of correction 180
to $\bZ(G)$ are %%
%% END of correction 180
multiples of the trivial character. (Let $\chi$ be
such a character and %%
%% Correction 181 [ ]
%% Annotated text: "χZpGq “ χp1qα<Caret> for some"
%% Comment: "" 
%% 
%% START of correction 181
let $\chi_{\bZ(G)}=\chi(1)\alpha$ for %%
%% END of correction 181
some
$\alpha\in\Irr(\bZ(G))$. %%
%% Correction 182 [ ]
%% Annotated text: "detpχp1qαq “ αχp1q<Caret>, which implies"
%% Comment: "" 
%% 
%% Correction 183 [ ]
%% Annotated text: "detpχqZpGq “ detpχZpGqq<Caret> “ detpχp1qαq"
%% Comment: "" 
%% 
%% START of corrections 182, 183
Then
$\mathbf{1}_{\bZ(G)}=\det(\chi)_{\bZ(G)}=\det(\chi_{\bZ(G)})=\det(\chi(1)\alpha)=\alpha^{\chi(1)}$,
which %%
%% END of corrections 182, 183
implies that the order of $\alpha$ divides
$(\chi(1),|\bZ(G)|)$.) This completes the proof.
\end{proof}

%%%%%%%%%%%%%%%%%%%%%%%
%%%%%%%%%%%%%%%%%%%%%%%
\section{A Proof of Theorem \ref{thm:main}}

We are now in the position to give a complete proof of our main result, as stated in the introduction.

A finite group admitting a Hall $\pi$-subgroup is often called an
\emph{$E_\pi$-group}. There is a large literature on the theory of
$E_\pi$-groups, including the results about simple groups we cited in
the previous section. We refer the reader to \cite{RV10} for the
latest results and relevant information. Note that the class of
$E_\pi$-groups is not closed under extensions and a subgroup of an
$E_\pi$-group might be no longer an $E_\pi$-group. Nevertheless, the
only fact we need is that every normal/subnormal subgroup and every
quotient of an $E_\pi$-group is an $E_\pi$-group. In fact, if $H$ is
a Hall $\pi$-subgroup of $G$ and $N\nor G$, then $H\cap N$ is a Hall
$\pi$-subgroup of $N$ %%
%% Correction 177 [ ]
%% Annotated text: "N and HN{N<Caret> is a"
%% Comment: "" 
%% 
%% START of correction 177
and $HN/N$ is %%
%% END of correction 177
%%
%% Correction 178 [ ]
%% Annotated text: "is a Hall<Caret> of G{N."
%% Comment: "" 
%% 
%% Correction 179 [ ]
%% Annotated text: "a Hall π<Caret>-subgroup of G{N."
%% Comment: "" 
%% 
%% START of corrections 178, 179
a Hall $\pi$-subgroup of %%
%% END of corrections 178, 179
$G/N$%%
%% Correction 174 [ ]
%% Annotated text: "Qpθq Ď Qpζpq<Caret> for some"
%% Comment: "" 
%% 
%% Correction 175 [ ]
%% Annotated text: "N IJ G<Caret> such that"
%% Comment: "" 
%% 
%% Correction 176 [ ]
%% Annotated text: "N| “ r<Caret> is an"
%% Comment: "" 
%% 
%% START of corrections 174, 175, 176
.


Another fact we need is that if $N\nor G$ such that $|G:N|=r$ is an
odd prime and $\theta\in \Irr(N)$ with $\QQ(\theta)\subseteq
\QQ(\zeta_p)$ for some prime $p\neq r$, %%
%% END of corrections 174, 175, 176
then there exists an
irreducible constituent $\chi$ of $\theta^G$ such that
$\QQ(\chi)\subseteq \QQ(\zeta_p)$. This follows from
\cite[Lemma~2.2]{GHSV21}, for instance.

We restate Theorem~\ref{thm:main} for the reader's convenience.

\begin{theorem}
Let $\pi$ be a set of odd primes and $G$ a finite group such that
$G$ has a nontrivial Hall $\pi$-subgroup. Then $G$ possesses a
nontrivial $\pi'$-degree irreducible character with values %%
%% Correction 173 [ ]
%% Annotated text: "values in Qpζpq<Caret> for some"
%% Comment: "" 
%% 
%% START of correction 173
in
$\QQ(\zeta_p)$ for %%
%% END of correction 173
some $p\in\pi$.
\end{theorem}

\begin{proof}
We argue by induction on $|G|$. %%
%% Correction 172 [ ]
%% Annotated text: "M Ÿ G<Caret> be a"
%% Comment: "" 
%% 
%% START of correction 172
Let $M \triangleleft\, G$ be %%
%% END of correction 172
a
normal subgroup such that $G/M$ is simple.

First assume %%
%% Correction 171 [ ]
%% Annotated text: "assume that G{M<Caret> is abelian"
%% Comment: "" 
%% 
%% START of correction 171
that $G/M$ is %%
%% END of correction 171
abelian and that $|G/M| \in \{2\} \cup
\pi$. Then the inflation to $G$ of any nontrivial linear character
of $G/M$ satisfies the required conditions. Now suppose %%
%% Correction 170 [ ]
%% Annotated text: "r :“ |G{M|<Caret> is an"
%% Comment: "" 
%% 
%% START of correction 170
that $r :=
|G/M|$ is %%
%% END of correction 170
an odd prime not belonging to $\pi$. In particular, we
have $\pi(M) \supseteq \pi\cap\pi(G)$.

By the induction hypothesis and the fact that $M$ also has a Hall
$\pi$-subgroup, there exists $\psi \in \Irr_{\pi'}(M)$ and %%
%% Correction 169 [ ]
%% Annotated text: "p P π<Caret> such that"
%% Comment: "" 
%% 
%% START of correction 169
some $p
\in \pi$ such %%
%% END of correction 169
that $\QQ(\psi) \subseteq \QQ(\zeta_p)$. Let $\chi \in
\Irr(G)$ lie over $\psi$. %%
%% Correction 165 [ ]
%% Annotated text: "řr i“1 ψi<Caret> is the"
%% Comment: "" 
%% 
%% Correction 166 [ ]
%% Annotated text: "“ ψ or<Caret> χM “"
%% Comment: "" 
%% 
%% Correction 167 [ ]
%% Annotated text: "χM “ ψ<Caret> or χM"
%% Comment: "" 
%% 
%% Correction 168 [ ]
%% Annotated text: "Corollary 6.19], either<Caret> χM “"
%% Comment: "" 
%% 
%% START of corrections 165, 166, 167, 168
By~\cite[Corollary~6.19]{Isaacs1}, either
$\chi_M = \psi$ or $\chi_M = \sum_{i=1}^r \psi_i$ is %%
%% END of corrections 165, 166, 167, 168
the sum of the
$G$-conjugates of $\psi$.

In the latter case, we may take $\chi = \psi^G$, with a note %%
%% Correction 161 [ ]
%% Annotated text: "χp1q “ rψp1q<Caret> is a"
%% Comment: "" 
%% 
%% START of correction 161
that
$\chi(1) = r \psi(1)$ is %%
%% END of correction 161
%%
%% Correction 162 [ ]
%% Annotated text: "is a π1-<Caret> number and"
%% Comment: "" 
%% 
%% Correction 163 [ ]
%% Annotated text: "is a π1-<Caret> number and"
%% Comment: "" 
%% 
%% Correction 164 [ ]
%% Annotated text: "is a π<Caret>1- number and"
%% Comment: "" 
%% 
%% START of corrections 162, 163, 164
a $\pi'$-number and %%
%% END of corrections 162, 163, 164
$\QQ(\chi) \subseteq
\QQ(\psi)\subseteq \QQ(\zeta_p)$. In the former case, every
irreducible character of $G$ lying over $\psi$ is an extension of
$\psi$, and hence has $\pi'$-degree. Moreover, as mentioned above,
one of these extensions has values in $\QQ(\zeta_p)$, as required.

Finally, suppose that $G/M$ is non-abelian. Set $\pi_1 := \pi(G/M)
\cap \pi$. (Note that $\pi_1$ could be empty.) Again, the group
$G/M$ has a Hall $\pi$-subgroup, which is also a Hall
$\pi_1$-subgroup. Theorem~\ref{thm:simple} then yields a character
$\chi \in \Irr_{\pi_1'}(G/M)$ such that $\QQ(\chi) \subseteq
\QQ(\zeta_p)$ for some $p \in \pi_1$. Since $\chi(1)$ divides
$|G/M|$, the character $\chi$ is also of $\pi'$-degree. The desired
character is obtained by inflating $\chi$ to $G$.
\end{proof}

We conclude with a remark that, in view of the above proof, the
conclusion of Theorem~\ref{thm:main} remains valid when $2 \in \pi$,
provided that the group $G$ is $\pi$-separable.

%%%%%%%%%%%%%%%%%%%%%%
%%%%%%%%%%%%%%%%%%%%%%%

\begin{thebibliography}{ABCDEF}

\bibitem[B-S70]{Borel70}
A. Borel, R. W. Carter, C.\,W. Curtis, N. Iwahori, T.\,A. Springer,
and R. Steinberg, \emph{Seminar on algebraic groups and related
finite groups}, Lect. Notes in Math. 131, Springer-Verlag, Berlin,
1970.

\bibitem[Car85]{Carter85}
R.\,W. Carter, \emph{Finite groups of Lie type. Conjugacy classes
and complex characters}, Pure Appl. Math. (N. Y.) Wiley-Intersci.
Publ. John Wiley \& Sons, Inc., New York, 1985.

\bibitem[Atl]{Atl1}
  J.\,H. Conway, R.\,T. Curtis, S.\,P. Norton, R.\,A. Parker, and R.\,A. Wilson,
 {\it Atlas of finite groups}, Clarendon Press, Oxford, 1985.

%\bibitem[DM91]{Digne-Michel}
%  F. Digne and J. Michel, {\it Representations of finite
%  groups of Lie type}, London Mathematical Society Student Texts \textbf{95}, $2020$.

\bibitem[FT63]{Thompson63}
W. Feit and J. Thompson, Solvability of groups of odd order,
\emph{Pacific J. Math.} \textbf{13} (1963), 775--1029.

%\bibitem[GAP]{GAP}
%{The GAP Group}, \emph{GAP Groups, Algorithms, and Programming,
%Version 4.11.0}, 2020. \url{http://www.gap-system.org}




\bibitem[GM20]{GeckMalle}
M. Geck and G. Malle, {\it The character theory of finite groups of
Lie type: a guided tour}, Cambridge studies in advanced mathematics
\textbf{187}, Cambridge University Press, Cambridge, 2020.

\bibitem[GHSV21]{GHSV21}
%%
%% Correction 193 [ ]
%% Annotated text: "[GHSV21] E. Giannel<Caret>li, N. N."
%% Comment: "" 
%% 
%% START of correction 193
E. Giannelli, N.%%
%% END of correction 193
\,N. Hung, %%
%% Correction 184 [ ]
%% Annotated text: "A. A. Schaeffe<Caret>r Fry, and"
%% Comment: "" 
%% 
%% Correction 185 [ ]
%% Annotated text: "A. A. Schaef<Caret>er Fry, and"
%% Comment: "" 
%% 
%% Correction 186 [ ]
%% Annotated text: "A. A. Schaef<Caret>er Fry, and"
%% Comment: "" 
%% 
%% Correction 187 [ ]
%% Annotated text: "A. A. Schaef<Caret>er Fry, and"
%% Comment: "" 
%% 
%% Correction 188 [ ]
%% Annotated text: "A. A. Schae<Caret>ffer Fry, and"
%% Comment: "" 
%% 
%% Correction 189 [ ]
%% Annotated text: "A. A. Scha<Caret>effer Fry, and"
%% Comment: "" 
%% 
%% Correction 190 [ ]
%% Annotated text: "A. A. Sch<Caret>aeffer Fry, and"
%% Comment: "" 
%% 
%% Correction 191 [ ]
%% Annotated text: "A. A. Sc<Caret>haeffer Fry, and"
%% Comment: "" 
%% 
%% START of corrections 184, 185, 186, 187, 188, 189, 190, 191
A.\,A. Schaeffer Fry, %%
%% END of corrections 184, 185, 186, 187, 188, 189, 190, 191
and %%
%% Correction 192 [ ]
%% Annotated text: "and C. Val<Caret>lejo, Characters of"
%% Comment: "" 
%% 
%% START of correction 192
C. Vallejo,
Characters %%
%% END of correction 192
of $\pi'$-degree and small cyclotomic fields, \emph{Ann.
Mat. Pura Appl.} \textbf{200} (2021), 1055--1073.

\bibitem[GLS98]{GLS98}
D. Gorenstein, R. Lyons, and R. Solomon, \emph{The classification of
the finite simple groups}, Number 3; Part I, Chapter A: Almost
simple $\mathcal{K}$-groups, Mathematical Surveys and Monographs,
40.3, American Mathematical Society, Providence, RI, 1998.

\bibitem[Gro86]{G86}
F. Gross, On a conjecture of Philip Hall, \emph{Proc. Lond. Math.
Soc. (3)} \textbf{52} (1986), no. 3, 464--494.

\bibitem[Hal28]{Hall1}
{P. Hall}, A note on soluble groups,
\emph{J. London. Math. Soc.} \textbf{3} (1928), 98--105.

\bibitem[Hal56]{Hall2}
{P. Hall}, Theorems like Sylow's,
\emph{Proc. London Math. Soc.} \textbf{6} (1956), 286--304.



\bibitem[Hun24]{Hung22}
{N.\,N. Hung}, The continuity of $p$-rationality and a lower bound
for $p'$-degree irreducible characters of finite groups,
\emph{Trans. Amer. Math. Soc.} \textbf{377} (2024), 323--344.

\bibitem[Isa06]{Isaacs1}
I.\,M. Isaacs, {\it Character theory of finite groups}, AMS Chelsea
Publishing, Providence, Rhode Island, 2006.

\bibitem[Mal07]{M07}
G. Malle, Height 0 characters of finite groups of Lie type,
\emph{Represent. Theory} \textbf{11} (2007), 192--220.

\bibitem[MT11]{malletesterman}
G.~Malle and D.~Testerman, \emph{Linear algebraic groups and finite
 groups of Lie type}, Cambridge studies in advanced mathematics
\textbf{133}, Cambridge University Press, Cambridge, 2011.


\bibitem[NT06]{Navarro-Tiep1}
G. Navarro and P.\,H. Tiep, Characters of $p'$-degree with
cyclotomic field of values, \emph{Proc. Amer. Math. Soc.}
\textbf{134} (2006), 2833--2837.

\bibitem[NT08]{Navarro-Tiep2}
 G. Navarro and P.\,H. Tiep, Rational irreducible
 characters and rational conjugacy classes in finite groups,
 \emph{Trans. Amer. Math. Soc.}
 \textbf{360} (2008), 2443--2465.

  \bibitem[Ras77]{Rasala}
R. Rasala, On the minimal degrees of characters of $\Sy_n$, \emph{J.
Algebra} \textbf{45} (1977), 132--181.

\bibitem[RV10]{RV10}
D.\,O. Revin and E.\,P. Vdovin, On the number of classes of
conjugate Hall subgroups in finite simple groups, \emph{J. Algebra}
\textbf{324} (2010), no. 12, 3614--3652.



\bibitem[Tho66]{Thompson}
J. Thompson, Hall subgroups of the symmetric groups,
\emph{J. Combinatorial Theory} \textbf{1} (1966), 271--279.

\end{thebibliography}
\end{document}
